\section{Consistent Evidence and Conservative Validation}
\label{sec:design}
A certificate must bind a value to the evidence that certifies it. A remembered fact is stored with its source, query/projection, value, support, adapter contract, provider epoch, evidence version and read interval. Entries without this schema, exceptions, missing evidence and unsupported contracts are \Unknown, not successful empty reads. Notifications are invalidation hints; silence is not a certificate.

\paragraph{Observation and re-anchoring.} Two protocols are possible. A cooperative provider can return value, support, epoch and sequence from one transaction. Alternatively, a reader obtains complete evidence $E_0$, executes the query and its support reads, and obtains $E_1$. Equality is sufficient only if the adapter establishes that every relevant committed change alters the evidence, the provider epoch cannot silently recur, and all reads obey the stated consistency model. Neither zero-miss trials nor equality of timestamps, repeated values or digests establishes those premises. A mismatching pair restarts the entire operation, for at most three attempts; exhaustion returns \Unknown. Refresh and re-anchoring use the same protocol, never a value from one attempt and a version from another. Algorithm~\ref{alg:validate} in Appendix~\ref{app:proofs} states the procedure for one source.

\paragraph{Validation's certified instant.} The protocol brackets the complete sequence of watermark, support, delta and query reads. A successful bracket certifies the query snapshot inside that stable interval. The log records its in-transaction timestamp and the enclosing read interval. \Changed denotes a witnessed value change; reset epochs and incomplete evidence yield \Unknown. A provider version change does not by itself prove that a particular fact changed.

\paragraph{Stock adapters.} We inspect complete Contacts version maps, MediaStore identity/generation and call-log row contents alongside the optimistic stamps. The stock workload uses joins, query membership and sometimes several cursors. The public Contacts columns give per-row versions and timestamps but no provider-wide sequence or epoch across the raw-contact, data and aggregate tables a query joins. Reading a row's value and version in one cursor would therefore not cover these paths: membership, deletions and joined rows need further queries, and a provider wiped and refilled reproduced its versions (\S\ref{sec:meas-gaps}). For MediaStore, a row's \texttt{generation\_modified} read with its stable ID covers at most that one row: it misses rows entering or leaving the query and deleted rows, and read apart from the value it was fooled in \PhysRaceMediaStampedStale race trials. Lacking an atomic epoch and query contract for these paths, strict stock certification stays disabled; direct refresh remains, with its read interval recorded.

\paragraph{Cooperative reference.} A separate Android provider process stores one keyed row and an epoch/sequence pair in SQLite. Every supported mutation and deletion increments the sequence in its transaction; managed reset or restore rotates the epoch in the same transaction. All writers must use this API. Because the query snapshot is already atomic, a stable bracket adds only retries. The reference is thus a one-row positive control: its \RevReference trials check that client and oracle follow the procedure through deletion, identifier reuse, reset and injected writes (guards in Appendix~\ref{app:proofs}). Its scope is one keyed row, one writer API and one provider process; multi-row queries, joins, sync adapters and stock providers lie outside it, and its bracket cost does not price a large provider.

\paragraph{Several sources.} Separately refreshed or certified inputs are each current at their own instant, and Android offers no joint snapshot across providers. Tasks correct when each input was current at its read, such as the drafts in \S\ref{sec:agent-tasks}, tolerate this; tasks that relate sources remain \Unknown. An effect that must not act on a changed input, such as calling a just-validated number, needs a provider-side conditional effect: re-validating just before it shortens, but no re-validation, fixed wait or fast callback delivery closes, the interval in which an overlapping write changes the fact (\S\ref{sec:eval-boundary}).
